%! AP Statistics — Unit 8, Lesson 8.3 — Student Packet Cover Sheet
\documentclass[10pt]{article}
\usepackage{apstats-article}
\usepackage{apstats-boxes}
\usepackage{tikz}
\usepackage{ltablex}
\keepXColumns

\begin{document}

\begin{tikzpicture}[remember picture, overlay]
  \fill[navy] (current page.north west) rectangle
    ([yshift=-2.2cm]current page.north east);
  \node[anchor=west, white, font=\large\bfseries] at
    ([xshift=1cm, yshift=-0.7cm]current page.north west)
    {\CourseName: \SchoolYear};
  \node[anchor=west, white, font=\normalsize] at
    ([xshift=1cm, yshift=-1.3cm]current page.north west)
    {Unit 8: Inference for Categorical Data: Chi-Square};
  \node[anchor=west, white, font=\normalsize\bfseries] at
    ([xshift=1cm, yshift=-1.9cm]current page.north west)
    {Lesson 8.3 \quad Carrying Out a Chi-Square Test for Goodness of Fit};
\end{tikzpicture}

\vspace{2.4cm}
\namedateperiod

\begin{learningtargetbox}
\begin{itemize}
    \item I can calculate the chi-square test statistic for a goodness-of-fit test
          using the formula $\chi^2 = \sum (O - E)^2 / E$.
    \item I can find and interpret the p-value for a chi-square goodness-of-fit test
          using a chi-square table.
    \item I can write a complete conclusion for a chi-square goodness-of-fit test
          in context, including comparison of $p$ to $\alpha$.
\end{itemize}
\end{learningtargetbox}

\begin{tocbox}
% A CONTENTS LIST ONLY: no Score column and no Total row. Whether a component counts is said in
% its description.
\renewcommand{\arraystretch}{1.5}
\begin{tabularx}{\linewidth}{@{} c l X @{}}
    \textbf{\#} & \textbf{Component} & \textbf{Description} \\
    \hline
    1 & Warm-Up        & Spiral review: hypotheses and expected counts (8.2) \\
    2 & Guided Notes   & Chi-square statistic, p-value, conclusion template \\
    3 & Group Activity & MLB birth month GOF test (full procedure) \\
    4 & Exit Ticket    & Candy color distribution: p-value and conclusion \\
    5 & Homework       & Three GOF problems (calculation + conceptual) \\
\end{tabularx}
\end{tocbox}

\end{document}
